\chapter{Chapter05 — Sharp Level One Inequality}
%%% generated by the blueprint pipeline from FABL.Chapter05.SharpLevelOneInequality
\section{Overview}

Declarations of the Lean module \texttt{FABL.Chapter05.SharpLevelOneInequality}.

\section{Declarations}

\begin{theorem}[Pointwise Gibbs inequality]
\lean{FABL.gibbs_pointwise}
\label{FABL.gibbs_pointwise}
\leanok
\difficulty{5}
For real numbers $q,y,Z$ with $q\ge 0$ and $Z>0$, one has $qy-q\log Z-q\log q\le
e^y/Z-q$, where the term $q\log q$ is interpreted using the usual continuous convention
at $q=0$.
\end{theorem}

\begin{theorem}[Entropy variational inequality]
\lean{FABL.expect_mul_sub_log_expect_exp_le_entropy}
\label{FABL.expect_mul_sub_log_expect_exp_le_entropy}
\leanok
\uses{FABL.gibbs_pointwise}
\difficulty{4}
Let $\iota$ be a finite nonempty set, and let $q,y:\iota\to\mathbb R$ satisfy $q(x)\ge
0$ for every $x$ and have uniform mean $1$. Then the uniform expectation of $q(x)y(x)$
minus the logarithm of the uniform expectation of $\exp(y(x))$ is at most the uniform
expectation of $q(x)\log q(x)$.
\end{theorem}

\begin{theorem}[Entropy bound from a pointwise upper bound]
\lean{FABL.expect_mul_log_le_log_bound}
\label{FABL.expect_mul_log_le_log_bound}
\leanok
\difficulty{4}
Let $\iota$ be a finite nonempty set, and let $q:\iota\to\mathbb R$ satisfy $0\le
q(x)\le M$ for every $x\in\iota$, with uniform mean $1$. Then the uniform expectation of
$q(x)\log q(x)$ is at most $\log M$.
\end{theorem}

\begin{theorem}[Exponential moment of a sign linear form]
\lean{FABL.expect_exp_linearForm_le}
\label{FABL.expect_exp_linearForm_le}
\leanok
\uses{FABL.Sign, FABL.expect_sign, FABL.linearForm, FABL.signValue}
\difficulty{4}
Let $a_1,\ldots,a_n\in\mathbb R$, and choose $x$ uniformly from $\{-1,1\}^n$. Then
$\mathbb E\exp\!\left(\sum_{i=1}^n a_i x_i\right)\leq \exp\!\left(\frac12\sum_{i=1}^n
a_i^2\right)$.
\end{theorem}

\begin{theorem}[Vanishing level-one Fourier weight]
\lean{FABL.sharpLevelOneInequality_eq_zero}
\label{FABL.sharpLevelOneInequality_eq_zero}
\leanok
\uses{FABL.fourierCoeff, FABL.fourierWeightAtLevel, FABL.fourierWeightAtLevel_one_eq_sum_singleton, FABL.mean}
\difficulty{3}
Let $n$ be a natural number and let $f$ be a real-valued function on the $n$-dimensional
Boolean cube $\\{-1,1\\}^n$ whose values lie in $[0,1]$. If the mean of $f$ on the cube
is zero, then its total Fourier weight at level one is zero.
\end{theorem}

\begin{theorem}[Sharp level one inequality]
\lean{FABL.sharpLevelOneInequality}
\statementsource{Analysis-of-Boolean-Functions-by-Ryan-ODonnell}{
  pdf-ffbe9c7b3306-p126-b002,
  pdf-ffbe9c7b3306-p126-b003,
  pdf-ffbe9c7b3306-p273-b005
}
\label{FABL.sharpLevelOneInequality}
\leanok
\uses{FABL.expect_exp_linearForm_le, FABL.expect_mul_linearForm_eq_singletonCoeffs, FABL.expect_mul_log_le_log_bound, FABL.expect_mul_sub_log_expect_exp_le_entropy, FABL.fourierCoeff, FABL.fourierWeightAtLevel, FABL.fourierWeightAtLevel_one_eq_sum_singleton, FABL.linearForm, FABL.mean, FABL.monomial}
\difficulty{6}
% dataset-pipeline: stub -- prose pending build_informal_statements.py
Statement pending regeneration.
\end{theorem}

\begin{definition}[Signed level-one counterexample]
\lean{FABL.sharpLevelOneSignedCounterexample}
\label{FABL.sharpLevelOneSignedCounterexample}
\leanok
\uses{FABL.linearForm}
On the two-dimensional sign cube, the signed linear form with both coefficients equal to
$1/2$ is the function $x\mapsto (x_1+x_2)/2$.
\end{definition}

\begin{theorem}[Range of a two-coordinate linear form]
\lean{FABL.sharpLevelOneSignedCounterexample_mem_Icc}
\label{FABL.sharpLevelOneSignedCounterexample_mem_Icc}
\leanok
\uses{FABL.linearForm, FABL.sharpLevelOneSignedCounterexample, FABL.signValue_eq_neg_one_or_one}
\difficulty{4}
For each pair $x=(x_1,x_2)\\in\{−1,1\}^2$, the value $\frac12(x_1+x_2)$ lies in the
interval $[-1,1]$.
\end{theorem}

\begin{theorem}[Mean absolute value of a two-coordinate form]
\lean{FABL.mean_abs_sharpLevelOneSignedCounterexample}
\label{FABL.mean_abs_sharpLevelOneSignedCounterexample}
\leanok
\uses{FABL.Sign, FABL.expect_sign, FABL.linearForm, FABL.mean, FABL.sharpLevelOneSignedCounterexample, FABL.signValue}
\difficulty{4}
On the uniform two-dimensional sign cube, the function $g(x_1,x_2)=(x_1+x_2)/2$ has mean
absolute value $1/2$.
\end{theorem}

\begin{theorem}[Singleton coefficient of a linear form]
\lean{FABL.fourierCoeff_linearForm_singleton}
\label{FABL.fourierCoeff_linearForm_singleton}
\leanok
\uses{FABL.expect_monomial_mul, FABL.fourierCoeff, FABL.linearForm, FABL.monomial, FABL.signValue}
\difficulty{5}
Let $n$ be a nonnegative integer, let $a$ assign a real number to each of the $n$
coordinates, and let $i$ be one of those coordinates. The Fourier coefficient of the
linear form $x\\mapsto\\sum_{r=0}^{n-1}a_r x_r$ at the singleton set $\\{i\\}$ is $a_i$.
\end{theorem}

\begin{theorem}[Level-One Weight of a Linear Function]
\lean{FABL.fourierWeightAtLevel_one_sharpLevelOneSignedCounterexample}
\label{FABL.fourierWeightAtLevel_one_sharpLevelOneSignedCounterexample}
\leanok
\uses{FABL.fourierCoeff_linearForm_singleton, FABL.fourierWeightAtLevel, FABL.fourierWeightAtLevel_one_eq_sum_singleton, FABL.sharpLevelOneSignedCounterexample}
\difficulty{3}
On the two-dimensional sign cube, let $g(x_1,x_2)=(x_1+x_2)/2$. The sum of the squared
Fourier coefficients of $g$ over singleton subsets is $1/2$.
\end{theorem}

\begin{theorem}[Failure of the level-one Fourier bound]
\lean{FABL.not_sharpLevelOneInequality_signed}
\label{FABL.not_sharpLevelOneInequality_signed}
\leanok
\uses{FABL.fourierWeightAtLevel, FABL.fourierWeightAtLevel_one_sharpLevelOneSignedCounterexample, FABL.mean, FABL.mean_abs_sharpLevelOneSignedCounterexample, FABL.sharpLevelOneSignedCounterexample, FABL.sharpLevelOneSignedCounterexample_mem_Icc}
\difficulty{4}
There exist a natural number $n$, a real-valued function $f$ on the $n$-dimensional sign
cube, and a real number $\alpha$ such that $|f(x)|\le 1$ for every input $x$, the mean
of $|f|$ is $\alpha$, and $0<\alpha\le 1/2$, while the total squared Fourier weight at
level one is strictly greater than $2\alpha^2\log(1/\alpha)$.
\end{theorem}
